### 1.1. Definitions
For an integer \(k\ge 2\), a positive integer \(n\) is called **\(k\)-full** if
\[
\forall\text{ primes }p:\ p\mid n\ \Rightarrow\ p^k\mid n.
\]
Thus:
* \(2\)-full \(\Leftrightarrow\) powerful/squareful.
* \(3\)-full is sometimes called **cubeful**.

### 1.2. Ambiguity/misstatement in the provided problem text
The problem text asks:

> “Are there any 2-full \(n\) such that \(n+1\) is 3-full?”

but then cites consecutive examples \((8,9)\) with “8 is 3-full and 9 is 2-full” and
\((12167,12168)\) with “\(12167=23^3\) and \(12168=2^3\cdot 3^2\cdot 13^2\)” (again: first is 3-full,
second is 2-full).

So the examples match the **reversed** condition.

I will therefore separate:

* **Literal statement:** find \(n\) such that \(n\) is 2-full and \(n+1\) is 3-full.
* **Corrected statement (minimal change consistent with examples):** find \(n\) such that \(n\) is
  3-full and \(n+1\) is 2-full; and ask whether the only such pairs are \((8,9)\) and
  \((12167,12168)\).

---
